\documentclass[11pt]{amsart}

\usepackage[T1]{fontenc}
\usepackage{lmodern}
\usepackage{microtype}
\usepackage{mathtools,amssymb,amsthm}
\usepackage[hidelinks]{hyperref}

\hypersetup{
  pdftitle={A 24-cell weak 4-resolving set of K11 x K11 refutes a surmised exact value}
}

\newtheorem{theorem}{Theorem}[section]
\newtheorem{lemma}[theorem]{Lemma}
\newtheorem{proposition}[theorem]{Proposition}
\newtheorem{corollary}[theorem]{Corollary}
\theoremstyle{remark}
\newtheorem*{remark}{Remark}

\DeclareMathOperator{\wdim}{wdim}

\title[A 24-cell weak 4-resolving set of $K_{11}\times K_{11}$]
{A 24-cell weak 4-resolving set of $K_{11}\times K_{11}$
 refutes a surmised exact value}

\date{}

\subjclass[2020]{05C12, 05C76, 05C35}
\keywords{weak $k$-metric dimension, weak $k$-resolving set, direct product,
tensor product, complete graph, counterexample}

\begin{document}

\begin{abstract}
For a connected graph $G$ and an integer $k\ge1$, a set $S\subseteq V(G)$ is
\emph{weak $k$-resolving} if
$\Delta_S(x,y):=\sum_{z\in S}\lvert d(x,z)-d(y,z)\rvert\ge k$ for every pair of
distinct vertices $x,y$, and $\wdim_k(G)$ is the least size of such a set.
Farhan, Kuziak and Yero prove that
$\wdim_4(K_n\times K_n)\le 2n+1+\lfloor n/4\rfloor$ for every $n\ge9$ and
wonder, in the first item of their concluding remarks, whether this upper bound
is the exact value. It is not. We exhibit a $24$-cell weak $4$-resolving set of
$K_{11}\times K_{11}$, prove the matching lower bound
$\wdim_4(K_n\times K_n)\ge 2n+\lceil 2n/11\rceil$ for every $n\ge4$, and
conclude
\[
  \wdim_4(K_n\times K_n)=2n+\lfloor n/4\rfloor=24,\,27,\,29,\,31
  \quad\text{for } n=11,12,13,14,
\]
one less than the surmised $25,28,30,32$ in each case, together with the
infinite family $\wdim_4(K_{11m}\times K_{11m})=24m$ for every $m\ge1$. The
upper bound of the source is true and is not contradicted.
\end{abstract}

\maketitle

\section{Introduction}

All graphs are finite, simple and connected. For vertices $x,y$ of $G$ and
$S\subseteq V(G)$ put
\[
  \Delta_S(x,y)=\sum_{z\in S}\bigl\lvert d(x,z)-d(y,z)\bigr\rvert .
\]
Following Peterin, Sedlar, \v{S}krekovski and Yero \cite{PSSY}, $S$ is a
\emph{weak $k$-resolving set} of $G$ if $\Delta_S(x,y)\ge k$ for every pair of
distinct $x,y\in V(G)$, and the \emph{weak $k$-metric dimension} $\wdim_k(G)$
is the smallest cardinality of a weak $k$-resolving set. The condition is
imposed on \emph{every} pair, including pairs of vertices that lie in $S$.

The \emph{direct} (tensor, categorical) product $G\times H$ has vertex set
$V(G)\times V(H)$, with $(g,h)$ adjacent to $(g',h')$ if and only if
$gg'\in E(G)$ \emph{and} $hh'\in E(H)$. Throughout, $n\ge3$ and we identify
$V(K_n\times K_n)$ with the cells of an $n\times n$ array,
$[n]\times[n]$, written $(\text{row},\text{column})$ and indexed from $1$. For
$n\ge3$ the graph $K_n\times K_n$ is connected of diameter $2$, and for
$u=(i,j)$, $v=(i',j')$,
\begin{equation}\label{eq:dist}
  d(u,v)=
  \begin{cases}
    0,&u=v,\\
    1,&i\ne i'\ \text{and}\ j\ne j',\\
    2,&\text{otherwise (exactly one coordinate agrees).}
  \end{cases}
\end{equation}
The Cartesian analogue is settled: $\wdim_4(K_n\,\square\,K_n)=2n$, a special
case of the treatment of Hamming graphs in \cite{FKKMY}.

\subsection*{The statement under test}

A.~Farhan, D.~Kuziak and I.~G. Yero, \emph{The weak $k$-metric dimension of the
direct product of complete graphs} \cite{FKY}, determine
$\wdim_k(K_n\times K_n)$ for every $n\ge4$ and every admissible $k$ except
$k=4$, where they prove an upper bound. In the notation of \cite{FKY},
Theorem~3.3 reads
\begin{equation}\label{eq:thm33}
  \wdim_4(K_n\times K_n)\le 2n+1+\Bigl\lfloor\frac n4\Bigr\rfloor
  \qquad\text{for every integer } n\ge9 ,
\end{equation}
and the first item of the concluding remarks of \cite{FKY}, introduced by
``the following open questions might be of interest'', reads verbatim
\begin{quote}\small
``Since the case $k=4$ is the only one in which an exact formula was not
deduced, it is clearly of interest to complete this case, and we wonder on
whether the exact value is precisely that shown as an upper bound in
Theorem~[3.3].''
\end{quote}
The bracketed number replaces a cross-reference of the source, which resolves to
Theorem~3.3, i.e.\ to \eqref{eq:thm33}. Resolved into a quantified statement,
what is surmised is
\begin{equation}\label{eq:surmise}
  \wdim_4(K_n\times K_n)=2n+1+\Bigl\lfloor\frac n4\Bigr\rfloor
  \qquad\text{for every integer } n\ge9 .
\end{equation}
It carries no statement number in \cite{FKY}, which is why it is quoted in full
rather than cited by number.

\subsection*{Results}

\begin{theorem}\label{thm:main}
The set of $24$ cells
\begin{equation}\label{eq:S}
\begin{aligned}
 S=\{&(1,1),(1,10),(2,1),(2,11),(3,2),(3,10),(4,2),(4,11),\\
     &(5,3),(5,10),(6,3),(6,11),(7,4),(7,5),(8,6),(8,7),\\
     &(9,8),(9,9),(10,4),(10,6),(10,8),(11,5),(11,7),(11,9)\}
\end{aligned}
\end{equation}
is a weak $4$-resolving set of $K_{11}\times K_{11}$. Hence
$\wdim_4(K_{11}\times K_{11})\le 24<25=2\cdot11+1+\lfloor11/4\rfloor$, and
\eqref{eq:surmise} is false.
\end{theorem}

\begin{theorem}\label{thm:lower}
For every integer $n\ge4$,
\[
  \wdim_4(K_n\times K_n)\ \ge\ 2n+\Bigl\lceil\frac{2n}{11}\Bigr\rceil .
\]
\end{theorem}

\begin{theorem}\label{thm:exact}
For $n=11,12,13,14$,
\[
  \wdim_4(K_n\times K_n)=2n+\Bigl\lfloor\frac n4\Bigr\rfloor
  =24,\ 27,\ 29,\ 31 ,
\]
which is exactly one less than the value $25,28,30,32$ surmised in
\eqref{eq:surmise}. Moreover
$\wdim_4(K_{11m}\times K_{11m})=24m$ for every integer $m\ge1$, whereas
\eqref{eq:surmise} gives $22m+1+\lfloor 11m/4\rfloor$ there; the gap
$1+\lfloor 11m/4\rfloor-2m$ grows linearly, like $3n/44$.
\end{theorem}

\section{The distance sums, and a two-line criterion}\label{sec:crit}

Let $S\subseteq[n]\times[n]$. Write $a_i=\lvert\{j:(i,j)\in S\}\rvert$ for the
$i$-th row degree and $b_j=\lvert\{i:(i,j)\in S\}\rvert$ for the $j$-th column
degree, so $\sum_ia_i=\sum_jb_j=\lvert S\rvert$. For distinct
$x=(i,j)$ and $y=(i',j')$ put $\alpha=\lvert\{x,y\}\cap S\rvert$, and when
$i\ne i'$ and $j\ne j'$ also
$\beta=\lvert S\cap\{(i,j'),(i',j)\}\rvert$ --- the number of the two remaining
corners of the rectangle spanned by $x$ and $y$ that lie in $S$.

\begin{lemma}\label{lem:delta}
Let $n\ge3$ and $x=(i,j)\ne y=(i',j')$ in $K_n\times K_n$. Then
\[
  \Delta_S(x,y)=
  \begin{cases}
    b_j+b_{j'}+\alpha, & i=i',\\[2pt]
    a_i+a_{i'}+\alpha, & j=j',\\[2pt]
    a_i+a_{i'}+b_j+b_{j'}-\alpha-2\beta, & i\ne i'\ \text{and}\ j\ne j'.
  \end{cases}
\]
\end{lemma}

This is Proposition~2.2 of \cite{FKY}; we include the proof to keep the note
self-contained.

\begin{proof}
Throughout, $z=(p,q)$ ranges over $S$ and distances are read off
\eqref{eq:dist}.

Suppose first $i=i'$ and $j\ne j'$. If $z=x$ then $d(x,z)=0$ while $z$ and $y$
share their row, so $d(y,z)=2$ and the term is $2$; symmetrically $z=y$
contributes $2$. For $z\notin\{x,y\}$: if $p=i$ then $z$ shares its row with
both $x$ and $y$, so $d(x,z)=d(y,z)=2$ and the term is $0$; if $p\ne i$ then
$d(x,z)=2$ exactly when $q=j$ and $d(y,z)=2$ exactly when $q=j'$, so the term
is $1$ precisely when $q\in\{j,j'\}$ and $0$ otherwise. The cells of $S$ with
$q=j$ and $p\ne i$ number $b_j-[x\in S]$, and those with $q=j'$ and $p\ne i$
number $b_{j'}-[y\in S]$. Altogether
$\Delta_S(x,y)=2\alpha+(b_j-[x\in S])+(b_{j'}-[y\in S])=b_j+b_{j'}+\alpha$.
Exchanging the roles of the two coordinates gives the case $j=j'$.

Now let $i\ne i'$ and $j\ne j'$, so $d(x,y)=1$. If $z=x$ the term is
$\lvert0-1\rvert=1$, and likewise for $z=y$; these contribute $\alpha$. For
$z\notin\{x,y\}$ we have $d(x,z),d(y,z)\in\{1,2\}$, and $d(x,z)=2$ exactly when
$z$ shares a coordinate with $x$, i.e.\ when $p=i$ or $q=j$. Hence the term is
$1$ precisely when $z$ shares a coordinate with exactly one of $x,y$. The cells
of $S$ sharing a coordinate with $x$ are $a_i+b_j-[x\in S]$ in number, of which
$x$ itself is one when $x\in S$, so $a_i+b_j-2[x\in S]$ of them lie in
$S\setminus\{x,y\}$ (note $y$ shares no coordinate with $x$); symmetrically for
$y$. A cell sharing a coordinate with both $x$ and $y$ satisfies
($p=i$ or $q=j$) and ($p=i'$ or $q=j'$); as $p$ cannot equal both $i$ and $i'$
and $q$ cannot equal both $j$ and $j'$, it is $(i,j')$ or $(i',j)$, so there
are $\beta$ of them and none is $x$ or $y$. By inclusion--exclusion the number
of $z\in S\setminus\{x,y\}$ sharing a coordinate with exactly one of $x,y$ is
\[
  \bigl(a_i+b_j-2[x\in S]\bigr)+\bigl(a_{i'}+b_{j'}-2[y\in S]\bigr)-2\beta
  =a_i+a_{i'}+b_j+b_{j'}-2\alpha-2\beta,
\]
and adding the contribution $\alpha$ of $z\in\{x,y\}$ gives the third case.
\end{proof}

Let $G_S$ be the bipartite graph with parts $\{r_1,\dots,r_n\}$ and
$\{c_1,\dots,c_n\}$ in which $r_ic_j$ is an edge if and only if $(i,j)\in S$;
thus $\lvert E(G_S)\rvert=\lvert S\rvert$, $\deg r_i=a_i$ and $\deg c_j=b_j$.
For a subgraph $F$ of $G_S$ write $T(F)$ for the sum of the $G_S$-degrees of
the vertices of $F$. Rectangles with three corners in $S$ are exactly the
paths with three edges in $G_S$, and rectangles with four corners in $S$ are
exactly the $4$-cycles of $G_S$.

\begin{lemma}\label{lem:crit}
Let $n\ge3$ and let $S\subseteq[n]\times[n]$ have $a_i\ge2$ and $b_j\ge2$ for
all $i,j$. Then $S$ is a weak $4$-resolving set of $K_n\times K_n$ if and only
if
\begin{enumerate}
\item[\textup{(P)}] $T(P)\ge9$ for every path $P$ with three edges in $G_S$,
      and
\item[\textup{(C)}] $T(C)\ge10$ for every $4$-cycle $C$ in $G_S$.
\end{enumerate}
\end{lemma}

\begin{proof}
Pairs with $i=i'$ or $j=j'$ are settled at once: by Lemma~\ref{lem:delta} their
$\Delta_S$ is at least $b_j+b_{j'}\ge4$, resp.\ $a_i+a_{i'}\ge4$, and they
involve neither (P) nor (C). So only pairs with $i\ne i'$ and $j\ne j'$ matter,
where Lemma~\ref{lem:delta} gives $\Delta_S(x,y)=T-\alpha-2\beta$ with
$T=a_i+a_{i'}+b_j+b_{j'}\ge8$, and the requirement $\Delta_S(x,y)\ge4$ reads
\begin{equation}\label{eq:need}
  \alpha+2\beta\le T-4 .
\end{equation}
Here $\alpha+\beta$ is the number of corners of the rectangle spanned by
$x,y$ that lie in $S$, and $\alpha$ counts those on the diagonal $\{x,y\}$.
If $\alpha+\beta\le2$ then $\alpha+2\beta\le4\le T-4$ and \eqref{eq:need}
holds. If $\alpha+\beta=3$ then either $(\alpha,\beta)=(2,1)$, when
$\alpha+2\beta=4\le T-4$ holds automatically, or $(\alpha,\beta)=(1,2)$, when
\eqref{eq:need} says exactly $T\ge9$. If $\alpha+\beta=4$ then
$(\alpha,\beta)=(2,2)$ and \eqref{eq:need} says exactly $T\ge10$.

Thus $S$ is weak $4$-resolving if and only if every rectangle with exactly
three corners in $S$ has $T\ge9$ and every rectangle with four corners in $S$
has $T\ge10$. The second condition is (C). Given (C), the first condition is
equivalent to (P): a three-edge path whose fourth rectangle corner also lies in
$S$ spans a $4$-cycle, for which (C) already gives $T\ge10\ge9$; and a
three-edge path whose fourth corner is outside $S$ is a rectangle with exactly
three corners in $S$, with the same $T$.
\end{proof}

\section{Proof of Theorem \ref{thm:main}}\label{sec:upper}

The $24$ cells of $S$ in \eqref{eq:S} are distinct, and the row and column
degree sequences are
\[
  (a_1,\dots,a_{11})=(b_1,\dots,b_{11})=(2,2,2,2,2,2,2,2,2,3,3),
\]
summing to $9\cdot2+2\cdot3=24$ both ways. In particular every degree is at
least $2$, so Lemma~\ref{lem:crit} applies. The row neighbourhoods in $G_S$ are
\begin{equation}\label{eq:nbhds}
\begin{gathered}
 N(r_1)=\{c_1,c_{10}\},\quad N(r_2)=\{c_1,c_{11}\},\quad
 N(r_3)=\{c_2,c_{10}\},\\
 N(r_4)=\{c_2,c_{11}\},\quad N(r_5)=\{c_3,c_{10}\},\quad
 N(r_6)=\{c_3,c_{11}\},\\
 N(r_7)=\{c_4,c_5\},\quad N(r_8)=\{c_6,c_7\},\quad
 N(r_9)=\{c_8,c_9\},\\
 N(r_{10})=\{c_4,c_6,c_8\},\quad N(r_{11})=\{c_5,c_7,c_9\}.
\end{gathered}
\end{equation}

\emph{Condition} (C) \emph{holds vacuously: $G_S$ has no $4$-cycle.} A
$4$-cycle is a pair of rows sharing two columns, so it suffices that any two of
the eleven sets in \eqref{eq:nbhds} meet in at most one element. Split the
columns into the block $B=\{c_1,c_2,c_3\}$, the block
$B'=\{c_{10},c_{11}\}$ and the block $M=\{c_4,\dots,c_9\}$. Each of
$N(r_1),\dots,N(r_6)$ consists of one column of $B$ and one column of $B'$, and
the six pairs so formed are the six distinct elements of $B\times B'$; two
distinct pairs agree in at most one coordinate. Each of
$N(r_7),\dots,N(r_{11})$ lies inside $M$, hence is disjoint from every
$N(r_i)$ with $i\le6$. Finally $N(r_7),N(r_8),N(r_9)$ are pairwise
disjoint, $N(r_{10})$ and $N(r_{11})$ are disjoint, and $N(r_{10})$ meets each
of $N(r_7),N(r_8),N(r_9)$ in exactly one column, as does $N(r_{11})$.

\emph{Condition} (P) \emph{holds because no three-edge path avoids a vertex of
degree~$3$.} The vertices of degree $3$ are $r_{10},r_{11},c_{10},c_{11}$; call
the remaining $18$ vertices \emph{light}. If a three-edge path $P$ contains a
degree-$3$ vertex then $T(P)\ge3+2+2+2=9$, as required. So it suffices to show
that the subgraph $L$ of $G_S$ induced by the light vertices contains no path
with three edges. The edges of $L$ are the twelve cells of $S$ lying in a light
row and a light column, namely
\[
\begin{gathered}
  (1,1),(2,1);\qquad (3,2),(4,2);\qquad (5,3),(6,3);\\
  (7,4),(7,5);\qquad (8,6),(8,7);\qquad (9,8),(9,9),
\end{gathered}
\]
and these form six vertex-disjoint two-edge paths, whose centres are
respectively $c_1$, $c_2$, $c_3$, $r_7$, $r_8$ and $r_9$. Every component of
$L$ therefore has
three vertices and two edges, and no three-edge path fits inside one.

By Lemma~\ref{lem:crit}, $S$ is a weak $4$-resolving set of
$K_{11}\times K_{11}$, so $\wdim_4(K_{11}\times K_{11})\le\lvert S\rvert=24$.
Since \eqref{eq:surmise} asserts the value $2\cdot11+1+\lfloor11/4\rfloor=25$
at $n=11$, it is false. \qed

\section{Proof of Theorem \ref{thm:lower}}\label{sec:lower}

Let $n\ge4$ and let $S$ be a weak $4$-resolving set of $K_n\times K_n$ with
$s=\lvert S\rvert$.

\emph{Step 1: any two row degrees sum to at least $4$, and likewise for
columns.} Fix $i\ne i'$. The cells of $S$ in rows $i$ and $i'$ occupy at most
$a_i+a_{i'}$ columns, so if $a_i+a_{i'}<n$ there is a column $j$ with
$(i,j)\notin S$ and $(i',j)\notin S$. Applying Lemma~\ref{lem:delta} to the
same-column pair $x=(i,j)$, $y=(i',j)$, which has $\alpha=0$, gives
$a_i+a_{i'}=\Delta_S(x,y)\ge4$. Hence $a_i+a_{i'}\ge\min(n,4)=4$, since
$n\ge4$. The column statement is symmetric.

\emph{Step 2: we may assume every degree is at least $2$.} Suppose
$a_{i_0}\le1$ for some $i_0$. By Step~1 every other row has degree at least
$4-a_{i_0}$, so $s\ge a_{i_0}+(n-1)(4-a_{i_0})$, which is $4n-4$ if
$a_{i_0}=0$ and $3n-2$ if $a_{i_0}=1$; in either case $s\ge3n-2$ because
$n\ge2$. Now $n\ge4$ gives $9n\ge33$, that is $\frac{2n}{11}+1\le n-2$, and
therefore $\lceil\frac{2n}{11}\rceil<\frac{2n}{11}+1\le n-2$, so
$s\ge3n-2=2n+(n-2)\ge2n+\lceil2n/11\rceil$ and we are done. The same argument
applies to a column of degree at most $1$. So assume from now on that
$a_i\ge2$ and $b_j\ge2$ for all $i,j$, and set $E=s-2n\ge0$.

\emph{Step 3: heavy vertices.} Call a vertex of $G_S$ \emph{heavy} if its
degree is at least $3$ and \emph{light} if its degree is exactly $2$, and let
$h$ be the number of heavy vertices. Since $G_S$ has $2n$ vertices and $s$
edges,
\[
  \sum_{v\in V(G_S)}(\deg v-2)=2s-4n=2E ,
\]
and every summand is $\ge0$ with the heavy ones $\ge1$; hence $h\le 2E$ and
\begin{equation}\label{eq:heavydeg}
  \sum_{v\ \mathrm{heavy}}\deg v=2h+2E .
\end{equation}

\emph{Step 4: the light part is a union of short paths.} Let $L$ be the
subgraph of $G_S$ induced by the light vertices. Every vertex of $L$ has
$L$-degree at most $2$. If $L$ contained a path with three edges, that path
would be a three-edge path of $G_S$ with $T=8$, contradicting
Lemma~\ref{lem:crit}(P). In particular $L$ contains no cycle: $G_S$ is
bipartite, so a shortest cycle of $L$ has length at least $4$ and any four
consecutive vertices of it form a three-edge path. Hence $L$ is a disjoint
union of paths, each with at most $3$ vertices.

\emph{Step 5: counting light--heavy edges.} Let $k$ be the number of
components of $L$. A component with $c\in\{1,2,3\}$ vertices has $c-1$ edges,
and the $G_S$-degrees of its vertices sum to $2c$, so exactly
$2c-2(c-1)=2$ edges leave it for heavy vertices. Summing over components, the
number of light--heavy edges of $G_S$ is exactly $2k$. Every such edge has a
heavy endpoint, so by \eqref{eq:heavydeg},
\[
  2k\ \le\ \sum_{v\ \mathrm{heavy}}\deg v\ =\ 2h+2E ,
  \qquad\text{i.e.}\qquad k\le h+E .
\]
On the other hand $L$ has $2n-h$ vertices distributed among $k$ components of
at most $3$ vertices each, so $k\ge(2n-h)/3$.

\emph{Step 6.} Combining, $2n-h\le3k\le3h+3E$, whence
\[
  2n\ \le\ 4h+3E\ \le\ 8E+3E\ =\ 11E .
\]
As $E$ is an integer, $E\ge\lceil2n/11\rceil$ and
$s=2n+E\ge2n+\lceil2n/11\rceil$. \qed

\begin{remark}
The hypothesis $n\ge4$ cannot be dropped: at $n=3$ Step~1 yields only
$a_i+a_{i'}\ge3$, and \cite{FKY} report
$\wdim_4(K_3\times K_3)=6<7=2\cdot3+\lceil6/11\rceil$.
\end{remark}

\section{Proof of Theorem \ref{thm:exact}}\label{sec:exact}

For $n=11$: Theorem~\ref{thm:main} gives $\wdim_4\le24$ and
Theorem~\ref{thm:lower} gives $\wdim_4\ge22+\lceil22/11\rceil=24$.

For $n=12,13,14$ the lower bounds are
$24+\lceil24/11\rceil=27$, $26+\lceil26/11\rceil=29$ and
$28+\lceil28/11\rceil=31$, and the matching upper bounds are witnessed by the
following sets, whose row and column degrees are all at least $2$ and which may
be checked against Lemma~\ref{lem:crit} or directly against
\eqref{eq:dist}:
\begin{align}
 S_{12}=\{&(1,11),(1,12),(2,5),(2,9),(3,7),(3,10),(3,12),(4,2),(4,4),(5,1),\notag\\
          &(5,3),(5,6),(6,6),(6,8),(7,2),(7,8),(7,11),(8,2),(8,9),(9,4),\notag\\
          &(9,9),(10,1),(10,10),(11,3),(11,7),(12,2),(12,5)\},\label{eq:S12}\\[4pt]
 S_{13}=\{&(1,6),(1,12),(2,7),(2,13),(3,1),(3,2),(3,5),(4,6),(4,10),(5,3),\notag\\
          &(5,5),(6,10),(6,12),(7,1),(7,7),(8,8),(8,11),(9,9),(9,10),(10,4),\notag\\
          &(10,8),(10,9),(11,3),(11,7),(11,11),(12,2),(12,4),(13,12),\notag\\
          &(13,13)\},\label{eq:S13}\\[4pt]
 S_{14}=\{&(1,12),(1,13),(2,3),(2,8),(3,1),(3,8),(4,6),(4,9),(4,14),(5,8),\notag\\
          &(5,10),(6,2),(6,4),(7,4),(7,10),(8,1),(8,2),(8,5),(9,7),(9,11),\notag\\
          &(9,12),(10,5),(10,11),(11,4),(11,9),(12,13),(12,14),(13,3),\notag\\
          &(13,13),(14,6),(14,7)\}.\label{eq:S14}
\end{align}
Here $\lvert S_{12}\rvert=27$, $\lvert S_{13}\rvert=29$ and
$\lvert S_{14}\rvert=31$. This proves the first assertion, and the surmised
values at these $n$ are $28,30,32$.

For the family, let $m\ge1$ and let
\[
  S(m)=\bigl\{\,(i+11t,\ j+11t)\ :\ 0\le t<m,\ (i,j)\in S\,\bigr\}
  \ \subseteq\ [11m]\times[11m]
\]
be the $m$-fold \emph{block-diagonal} copy of \eqref{eq:S}, of size
$24m$. Its bipartite graph $G_{S(m)}$ is the disjoint union of $m$ copies of
$G_S$, so all its degrees are $2$ or $3$, it has no $4$-cycle, and every
three-edge path lies inside one block and hence meets a vertex of degree $3$ by
the argument of Section~\ref{sec:upper}. Lemma~\ref{lem:crit} therefore gives
$\wdim_4(K_{11m}\times K_{11m})\le24m$. Theorem~\ref{thm:lower} with
$n=11m$ gives the reverse inequality,
\[
  \wdim_4(K_{11m}\times K_{11m})\ \ge\ 22m+\Bigl\lceil\frac{22m}{11}\Bigr\rceil
  \ =\ 22m+2m\ =\ 24m ,
\]
so equality holds. Finally $24m<22m+1+\lfloor11m/4\rfloor$ for every $m\ge1$,
since $1+\lfloor11m/4\rfloor-2m\ge1+\frac{11m-3}{4}-2m=\frac{3m+1}{4}>0$. \qed

\section{Scope}\label{sec:status}

Inequality \eqref{eq:thm33} is a theorem of \cite{FKY} and is not contradicted
here: $24\le25$, $27\le28$, $29\le30$ and $31\le32$. What is refuted is only
the surmised \emph{equality} \eqref{eq:surmise}, so this note is not an
erratum.

Theorem~\ref{thm:lower} gives $\wdim_4(K_9\times K_9)\ge20$ and
$\wdim_4(K_{10}\times K_{10})\ge22$, against the surmised $21$ and $23$;
combined with \eqref{eq:thm33} this leaves
$\wdim_4(K_9\times K_9)\in\{20,21\}$ and
$\wdim_4(K_{10}\times K_{10})\in\{22,23\}$. We settle neither, so
\eqref{eq:surmise} may well be correct at the two smallest members of its own
range. Theorem~\ref{thm:main} nevertheless refutes \eqref{eq:surmise}, which is
universally quantified over $n\ge9$.

Theorem~\ref{thm:exact} determines $\wdim_4(K_n\times K_n)$ at $n=11,12,13,14$
and at every multiple of $11$. For every other $n\ge9$ we have only
$2n+\lceil2n/11\rceil\le\wdim_4(K_n\times K_n)\le2n+1+\lfloor n/4\rfloor$, and
we do not know which end is attained. The lower bounds above are the hand proof
of Section~\ref{sec:lower}; no enumeration of candidate sets was carried out at
any $n$, and no claim of minimality rests on a solver.

\section{Verification}

Every object needed above is printed in this note, and the arguments of
Sections~\ref{sec:crit}--\ref{sec:exact} use nothing else. The accompanying
program \texttt{verify.py} reads the cell lists \eqref{eq:S},
\eqref{eq:S12}--\eqref{eq:S14} exactly as printed and uses the Python standard
library only. It checks \eqref{eq:dist} against breadth-first distances in the
explicitly constructed product for $3\le n\le7$; it checks
Lemma~\ref{lem:delta} against raw distance sums over all pairs; it recomputes
$\min\Delta_S$ over all pairs for each printed set and for the block-diagonal
sets $S(1),S(2),S(3)$, obtaining $\min\Delta_S=4$ in each case; and it verifies
the degree sequences, the absence of a $4$-cycle in $G_S$ and the
six-component structure of $L$. Of Theorem~\ref{thm:lower} it checks only the
arithmetic --- the chain $2n\le4h+3E\le11E$ for $4\le n\le300$ --- and not the
combinatorial Steps~1, 4 and 5, which are hand proofs; likewise the
block-diagonal family is machine-checked only at small $m$, the statement for
all $m\ge1$ being the hand proof of Section~\ref{sec:exact}. The recorded run
\texttt{verify.output.txt} reports $109$ passing checks and states in its own
output what it does not cover.

\begin{thebibliography}{9}

\bibitem{FKY}
A.~Farhan, D.~Kuziak, and I.~G. Yero,
\emph{The weak $k$-metric dimension of the direct product of complete graphs},
arXiv:2605.22307v1, 21 May 2026.
\href{https://arxiv.org/abs/2605.22307}{arXiv:2605.22307}.
Theorem and section numbers are quoted from v1, the only version at the time of
writing.

\bibitem{FKKMY}
J.~A. Fern\'andez, S.~Klav\v{z}ar, D.~Kuziak, M.~Mu\~noz-M\'arquez, and
I.~G. Yero,
\emph{On the weak $k$-metric dimension of Hamming graphs},
Discrete Optim. \textbf{60} (2026), 100945;
\href{https://arxiv.org/abs/2505.19642}{arXiv:2505.19642}.

\bibitem{PSSY}
I.~Peterin, J.~Sedlar, R.~\v{S}krekovski, and I.~G. Yero,
\emph{Resolving vertices of graphs with differences},
Comput. Appl. Math. \textbf{43} (2024), article~275;
\href{https://arxiv.org/abs/2309.00922}{arXiv:2309.00922}.
Introduces $\wdim_k$.

\end{thebibliography}

\end{document}
